\documentclass[../../main.tex]{subfiles}
\graphicspath{{../../image/}{../../figure/}}
\begin{document}

\section{黏性消失法}\label{section:ch06-黏性消失法}

此处的目的是把标量双曲守恒律\eqref{align:ch06-6-1-1}的可容许弱解构造为如下抛物方程组的解当 $\mu\downarrow 0$ 时的极限：
\begin{align}
\partial_{t}u(x,t)+\operatorname{div}G(u(x,t))=\mu\Delta u(x,t),\qquad x\in\mathbb{R}^{m},\ t\in[0,\infty),
\label{align:ch06-6-3-1}
\end{align}
其中 $\Delta$ 表示关于空间变量的 Laplace 算子，即 $\Delta=\sum_{\alpha=1}^{m}\partial_{\alpha}^{2}$，$\mu$ 为正参数。

该途径的动机已在\cref{section:ch04-消失黏性方法}给出。

注意\eqref{align:ch06-6-3-1}不必与任何具体物理模型相联系，因而项 $\mu\Delta u$ 应视为“人工黏性”。

由于\eqref{align:ch06-6-3-1}是抛物型的，初值问题\eqref{align:ch06-6-3-1}、\eqref{align:ch06-6-1-2}总有唯一解；即使初值 $u_{0}$ 仅属于 L$^{\infty}$，该解对 $t>0$ 也是光滑的（设 $G$ 正则）。例如，若导数 $G^{\prime}$ Hölder 连续，则\eqref{align:ch06-6-3-1}、\eqref{align:ch06-6-1-2}的解 $u$ 在 $\mathbb{R}^{m}\times(0,\infty)$ 上关于 $t$ 连续可微、关于空间变量两次连续可微。

把\eqref{align:ch06-6-1-1}、\eqref{align:ch06-6-1-2}的“相关”解视为\eqref{align:ch06-6-3-1}、\eqref{align:ch06-6-1-2}的解当 $\mu\downarrow 0$ 时的极限，这一前提为\cref{definition:ch06-6-2-1}所假定的可容许弱解概念提供了第一个依据：

\begin{theorem}\label{theorem:ch06-6-3-1}
设 $u_{\mu}$ 表示\eqref{align:ch06-6-3-1}、\eqref{align:ch06-6-1-2}的解。假设存在序列 $\{\mu_{k}\}$，$\mu_{k}\downarrow 0$（$k\to\infty$），使 $\{u_{\mu_{k}}\}$ 在 $\mathbb{R}^{m}\times[0,\infty)$ 上一致有界且几乎处处收敛于某函数 $u$。则 $u$ 是\eqref{align:ch06-6-1-1}、\eqref{align:ch06-6-1-2}在 $\mathbb{R}^{m}\times[0,\infty)$ 上的可容许弱解。
\end{theorem}

\begin{proof}
任取光滑凸熵函数 $\eta$，相应熵流 $Q$ 由\eqref{align:ch06-6-2-1}确定。以 $\eta^{\prime}(u_{\mu}(x,t))$ 乘\eqref{align:ch06-6-3-1}并利用\eqref{align:ch06-6-2-1}，得
\begin{align}
\partial_{t}\eta(u_{\mu})+\operatorname{div}Q(u_{\mu})=\mu\Delta\eta(u_{\mu})-\mu\eta^{\prime\prime}(u_{\mu})|\nabla u_{\mu}|^{2}.
\label{align:ch06-6-3-2}
\end{align}
以任一光滑非负、在 $\mathbb{R}^{m}\times[0,\infty)$ 上具紧支集的试验函数 $\psi$ 乘\eqref{align:ch06-6-3-2}，在 $\mathbb{R}^{m}\times[0,\infty)$ 上积分并分部积分。注意\eqref{align:ch06-6-3-2}右端最后一项非负，得不等式
\begin{align}
&\int_{0}^{\infty}\int_{\mathbb{R}^{m}}\Bigl{[}\partial_{t}\psi\,\eta(u_{\mu})+\sum_{\alpha=1}^{m}\partial_{\alpha}\psi\,Q_{\alpha}(u_{\mu})\Bigr{]}\,\mathrm{d}x\,\mathrm{d}t+\int_{\mathbb{R}^{m}}\psi(x,0)\eta(u_{0}(x))\,\mathrm{d}x\notag\\
&\qquad\geqslant-\mu\int_{0}^{\infty}\int_{\mathbb{R}^{m}}\Delta\psi\,\eta(u_{\mu})\,\mathrm{d}x\,\mathrm{d}t.
\label{align:ch06-6-3-3}
\end{align}
在\eqref{align:ch06-6-3-3}中取 $\mu=\mu_{k}$ 并令 $k\to\infty$，可知 $\{u_{\mu_{k}}\}$ 的极限 $u$ 对所有光滑凸熵 $\eta$ 与所有光滑非负试验函数 $\psi$ 满足\eqref{align:ch06-6-2-3}。由完备性可知，即使 $\eta$、$\psi$ 仅 Lipschitz 连续，\eqref{align:ch06-6-2-3}仍成立。证明完成。
\end{proof}

将\cref{theorem:ch06-6-2-3}与下述定理比较，即可明白\eqref{align:ch06-6-1-1}与\eqref{align:ch06-6-3-1}是完全相配的：

\begin{theorem}\label{theorem:ch06-6-3-2}
设 $u_{\mu}$ 与 $\bar{u}_{\mu}$ 是\eqref{align:ch06-6-3-1}以 $u_{0}$、$\bar{u}_{0}\in L^{1}(\mathbb{R}^{m})$ 为相应初值的解，且取值于紧区间 $[a,b]$。则对任意 $t>0$，
\begin{align}
\int_{\mathbb{R}^{m}}[u_{\mu}(x,t)-\bar{u}_{\mu}(x,t)]^{+}\,\mathrm{d}x\leqslant\int_{\mathbb{R}^{m}}[u_{0}(x)-\bar{u}_{0}(x)]^{+}\,\mathrm{d}x,
\label{align:ch06-6-3-4}
\end{align}
\begin{align}
\|u_{\mu}(\cdot,t)-\bar{u}_{\mu}(\cdot,t)\|_{L^{1}(\mathbb{R}^{m})}\leqslant\|u_{0}(\cdot)-\bar{u}_{0}(\cdot)\|_{L^{1}(\mathbb{R}^{m})}.
\label{align:ch06-6-3-5}
\end{align}
此外，若
\begin{align}
u_{0}(x)\leqslant\bar{u}_{0}(x),\qquad\text{a.e.\ 在 }\mathbb{R}^{m}\text{ 上},
\label{align:ch06-6-3-6}
\end{align}
则
\begin{align}
u_{\mu}(x,t)\leqslant\bar{u}_{\mu}(x,t),\qquad\text{在 }\mathbb{R}^{m}\times(0,\infty)\text{ 上}.
\label{align:ch06-6-3-7}
\end{align}
特别地，$u_{\mu}$ 与 $\bar{u}_{\mu}$ 的值域均含于 $[a,b]$。
\end{theorem}

\begin{proof}
为简化记号，省略下标 $\mu$，把 $u_{\mu}$、$\bar{u}_{\mu}$ 记作 $u$、$\bar{u}$。由抛物方程的标准理论可知，当 $u_{0}(\cdot)$、$\bar{u}_{0}(\cdot)\in L^{1}(\mathbb{R}^{m})\cap L^{\infty}(\mathbb{R}^{m})$ 时，$u(\cdot,t)$、$\bar{u}(\cdot,t)$ 及其任意阶空间导数也属于 $L^{1}(\mathbb{R}^{m})\cap L^{\infty}(\mathbb{R}^{m})$，且在 $(0,\infty)$ 的紧子集上关于 $t$ 一致有界。

对 $\varepsilon>0$，在 $\mathbb{R}$ 上定义函数
\begin{align}
\eta_{\varepsilon}(w)=
\begin{cases}
0, & -\infty<w\leqslant 0,\\[2pt]
\dfrac{w^{2}}{4\varepsilon}, & 0<w\leqslant 2\varepsilon,\\[6pt]
w-\varepsilon, & 2\varepsilon<w<\infty.
\end{cases}
\label{align:ch06-6-3-8}
\end{align}
由于 $u$ 与 $\bar{u}$ 都满足\eqref{align:ch06-6-3-1}，容易验证方程
\begin{align}
&\partial_{t}\eta_{\varepsilon}(u-\bar{u})+\sum_{\alpha=1}^{m}\partial_{\alpha}\{\eta_{\varepsilon}^{\prime}(u-\bar{u})[G_{\alpha}(u)-G_{\alpha}(\bar{u})]\}\notag\\
&\qquad-\sum_{\alpha=1}^{m}\eta_{\varepsilon}^{\prime\prime}(u-\bar{u})[G_{\alpha}(u)-G_{\alpha}(\bar{u})]\partial_{\alpha}(u-\bar{u})\notag\\
&\qquad\qquad=\mu\Delta\eta_{\varepsilon}(u-\bar{u})-\mu\eta_{\varepsilon}^{\prime\prime}(u-\bar{u})|\nabla(u-\bar{u})|^{2}.
\label{align:ch06-6-3-9}
\end{align}
固定 $0<s<t<\infty$，在 $\mathbb{R}^{m}\times(s,t)$ 上积分\eqref{align:ch06-6-3-9}。注意\eqref{align:ch06-6-3-9}右端最后一项非负，得不等式
\begin{align}
&\int_{\mathbb{R}^{m}}\eta_{\varepsilon}(u(x,t)-\bar{u}(x,t))\,\mathrm{d}x-\int_{\mathbb{R}^{m}}\eta_{\varepsilon}(u(x,s)-\bar{u}(x,s))\,\mathrm{d}x\notag\\
&\qquad\leqslant\sum_{\alpha=1}^{m}\int_{s}^{t}\int_{\mathbb{R}^{m}}\eta_{\varepsilon}^{\prime\prime}(u-\bar{u})[G_{\alpha}(u)-G_{\alpha}(\bar{u})]\partial_{\alpha}(u-\bar{u})\,\mathrm{d}x\,\mathrm{d}\tau.
\label{align:ch06-6-3-10}
\end{align}
注意 $\eta_{\varepsilon}^{\prime\prime}(u-\bar{u})[G_{\alpha}(u)-G_{\alpha}(\bar{u})]$ 关于 $\varepsilon>0$ 一致有界。又显然，当 $\varepsilon\downarrow 0$ 时 $\eta_{\varepsilon}(u(x,t)-\bar{u}(x,t))$ 逐点收敛于 $[u(x,t)-\bar{u}(x,t)]^{+}$，而 $\eta_{\varepsilon}^{\prime\prime}(u(x,t)-\bar{u}(x,t))[G_{\alpha}(u(x,t))-G_{\alpha}(\bar{u}(x,t))]$ 逐点收敛于零。因此由\eqref{align:ch06-6-3-10}与 Lebesgue 控制收敛定理，
\begin{align}
\int_{\mathbb{R}^{m}}[u(x,t)-\bar{u}(x,t)]^{+}\,\mathrm{d}x-\int_{\mathbb{R}^{m}}[u(x,s)-\bar{u}(x,s)]^{+}\,\mathrm{d}x\leqslant 0,
\label{align:ch06-6-3-11}
\end{align}
令 $s\downarrow 0$ 即得\eqref{align:ch06-6-3-4}。

在\eqref{align:ch06-6-3-4}中交换 $u$ 与 $\bar{u}$ 的地位，得类似不等式，与\eqref{align:ch06-6-3-4}相加即得\eqref{align:ch06-6-3-5}。

由\eqref{align:ch06-6-3-4}，\eqref{align:ch06-6-3-6}显然蕴含\eqref{align:ch06-6-3-7}。特别地，先取 $\bar{u}_{0}(x)\equiv b$、再取 $u_{0}(x)\equiv a$ 应用此单调性，得 $u(x,t)\leqslant b$、$\bar{u}(x,t)\geqslant a$。交换 $u$ 与 $\bar{u}$ 的地位，可知两解的值域均含于 $[a,b]$。证明完成。
\end{proof}

估计\eqref{align:ch06-6-3-5}可用于估计初值属于 $L^{\infty}(\mathbb{R}^{m})\cap L^{1}(\mathbb{R}^{m})$ 时\eqref{align:ch06-6-3-1}解的平均连续模。

\begin{lemma}\label{lemma:ch06-6-3-3}
设 $u_{\mu}$ 为\eqref{align:ch06-6-3-1}、\eqref{align:ch06-6-1-2}的解，其中 $u_{0}\in L^{1}(\mathbb{R}^{m})$ 且取值于紧区间 $[a,b]$。特别地，
\begin{align}
\int_{\mathbb{R}^{m}}|u_{0}(x+y)-u_{0}(x)|\,\mathrm{d}x\leqslant\omega(|y|),\qquad y\in\mathbb{R}^{m},
\label{align:ch06-6-3-12}
\end{align}
其中 $\omega$ 是 $[0,\infty)$ 上某非减函数，$\omega(r)\downarrow 0$（$r\downarrow 0$）。存在只依赖于 $[a,b]$ 的常数 $c$，使对任意 $t>0$，
\begin{align}
\int_{\mathbb{R}^{m}}|u_{\mu}(x+y,t)-u_{\mu}(x,t)|\,\mathrm{d}x\leqslant\omega(|y|),\qquad y\in\mathbb{R}^{m},
\label{align:ch06-6-3-13}
\end{align}
\begin{align}
\int_{\mathbb{R}^{m}}|u_{\mu}(x,t+h)-u_{\mu}(x,t)|\,\mathrm{d}x\leqslant c(h^{2/3}+\mu h^{1/3})\|u_{0}\|_{L^{1}(\mathbb{R}^{m})}+2\omega(h^{1/3}),\qquad h>0.
\label{align:ch06-6-3-14}
\end{align}
\end{lemma}

\begin{proof}
固定 $t>0$。对任意 $y\in\mathbb{R}^{m}$，函数 $\bar{u}_{\mu}(x,t)=u_{\mu}(x+y,t)$ 是\eqref{align:ch06-6-3-1}以 $\bar{u}_{0}(x)=u_{0}(x+y)$ 为初值的解。应用\eqref{align:ch06-6-3-5}，得
\begin{align}
\int_{\mathbb{R}^{m}}|u_{\mu}(x+y,t)-u_{\mu}(x,t)|\,\mathrm{d}x\leqslant\int_{\mathbb{R}^{m}}|u_{0}(x+y)-u_{0}(x)|\,\mathrm{d}x,
\label{align:ch06-6-3-15}
\end{align}
由此得\eqref{align:ch06-6-3-13}。

现固定 $h>0$。通过减去 $G(0)$ 归一化 $G$，因而以后可不失一般性地设 $G(0)=0$。以 $\mathbb{R}^{m}$ 上有界光滑函数 $\phi$ 乘\eqref{align:ch06-6-3-1}，在带形区域 $\mathbb{R}^{m}\times(t,t+h)$ 上积分。分部积分得
\begin{align}
\int_{\mathbb{R}^{m}}\phi(x)[u_{\mu}(x,t+h)-u_{\mu}(x,t)]\,\mathrm{d}x
=\int_{t}^{t+h}\int_{\mathbb{R}^{m}}\Bigl{\{}\sum_{\alpha=1}^{m}\partial_{\alpha}\phi(x)\,G_{\alpha}(u_{\mu}(x,\tau))+\mu\Delta\phi(x)\,u_{\mu}(x,\tau)\Bigr{\}}\,\mathrm{d}x\,\mathrm{d}\tau.
\label{align:ch06-6-3-16}
\end{align}
令
\begin{align}
v(x)=u_{\mu}(x,t+h)-u_{\mu}(x,t).
\label{align:ch06-6-3-17}
\end{align}
形式上，在\eqref{align:ch06-6-3-16}中插入 $\phi(x)=\operatorname{sgn}v(x)$ 即可建立\eqref{align:ch06-6-3-14}。但由于函数 $\operatorname{sgn}$ 不连续，须先用磨光子磨光它：取 $\mathbb{R}$ 上光滑非负、支集含于 $[-m^{-1/2},m^{-1/2}]$、总质量为一的函数 $\rho$（见\eqref{align:ch06-6-2-15}），令
\begin{align}
\phi(x)=\int_{\mathbb{R}^{m}}h^{-m/3}\prod_{\beta=1}^{m}\rho\Bigl{(}\frac{x_{\beta}-z_{\beta}}{h^{1/3}}\Bigr{)}\operatorname{sgn}v(z)\,\mathrm{d}z.
\label{align:ch06-6-3-18}
\end{align}
注意 $|\partial_{\alpha}\phi|\leqslant c_{1}h^{-1/3}$，$|\Delta\phi|\leqslant c_{2}h^{-2/3}$。此外，由\eqref{align:ch06-6-3-5}（取 $\bar{u}\equiv 0$），$\|u(\cdot,\tau)\|_{L^{1}(\mathbb{R}^{m})}\leqslant\|u_{0}(\cdot)\|_{L^{1}(\mathbb{R}^{m})}$。因此\eqref{align:ch06-6-3-16}蕴含
\begin{align}
\int_{\mathbb{R}^{m}}\phi(x)v(x)\,\mathrm{d}x\leqslant c(h^{2/3}+\mu h^{1/3})\|u_{0}\|_{L^{1}(\mathbb{R}^{m})},
\label{align:ch06-6-3-19}
\end{align}
其中 $c$ 只依赖于 $[a,b]$。另一方面，注意到
\begin{align}
|v(x)|-v(x)\operatorname{sgn}v(z)=|v(x)|-|v(z)|+[v(z)-v(x)]\operatorname{sgn}v(z)\leqslant 2|v(x)-v(z)|,
\label{align:ch06-6-3-20}
\end{align}
由\eqref{align:ch06-6-3-18}得
\begin{align}
|v(x)|-\phi(x)v(x)
&=\int_{\mathbb{R}^{m}}h^{-m/3}\prod_{\beta=1}^{m}\rho\Bigl{(}\frac{x_{\beta}-z_{\beta}}{h^{1/3}}\Bigr{)}[|v(x)|-v(x)\operatorname{sgn}v(z)]\,\mathrm{d}z\notag\\
&\leqslant 2\int_{|\xi|<1}\prod_{\beta=1}^{m}\rho(\xi_{\beta})\,|v(x)-v(x-h^{1/3}\xi)|\,\mathrm{d}\xi.
\label{align:ch06-6-3-21}
\end{align}
综合\eqref{align:ch06-6-3-17}、\eqref{align:ch06-6-3-21}、\eqref{align:ch06-6-3-19}与\eqref{align:ch06-6-3-13}，即得\eqref{align:ch06-6-3-14}。证明完成。
\end{proof}

至此已为给出\cref{theorem:ch06-6-2-2}的证明铺好基础。

\begin{proof}[\cref{theorem:ch06-6-2-2}的证明]
先设 $u_{0}\in L^{\infty}(\mathbb{R}^{m})\cap L^{1}(\mathbb{R}^{m})$。令 $u_{\mu}$ 表示\eqref{align:ch06-6-3-1}、\eqref{align:ch06-6-1-2}在 $0<\mu<1$ 时的解。由\eqref{align:ch06-6-3-14}，族 $\{u_{\mu}(\cdot,t)\}$（视为 $C^{0}([0,\infty);L^{1}(\mathbb{R}^{m}))$ 中的族）一致等度连续。此外，\eqref{align:ch06-6-3-13}蕴含：对任一固定 $t\in[0,\infty)$，$\{u_{\mu}(\cdot,t)\}$ 含于 $L^{1}_{\mathrm{loc}}(\mathbb{R}^{m})$ 的一个紧集中。于是由 Ascoli 定理，对任一序列 $\{\mu_{k}\}$（$\mu_{k}\to 0$，$k\to\infty$），存在子列（仍记为 $\{\mu_{k}\}$）及函数 $u\in C^{0}([0,\infty);L^{1}_{\mathrm{loc}}(\mathbb{R}^{m}))$，使 $\{u_{\mu_{k}}(\cdot,t)\}$ 在 $L^{1}_{\mathrm{loc}}(\mathbb{R}^{m})$ 中收敛于 $u(\cdot,t)$，且在 $[0,\infty)$ 的任一紧子集上关于 $t$ 一致。必要时再取子列（仍记为 $\{\mu_{k}\}$），可知 $\{u_{\mu_{k}}\}$ 在 $\mathbb{R}^{m}\times[0,\infty)$ 上一致有界且几乎处处收敛于 $u$；于是由\cref{theorem:ch06-6-3-1}，$u$ 是\eqref{align:ch06-6-1-1}、\eqref{align:ch06-6-1-2}的可容许弱解。由于可容许解唯一（\cref{corollary:ch06-6-2-4}），整个族 $\{u_{\mu}\}$ 当 $\mu\downarrow 0$ 时必收敛于 $u$。此外，由\eqref{align:ch06-6-3-14}及 $L^{1}(\mathbb{R}^{m})$ 中的弱下半连续性，$u(\cdot,t)\in C^{0}([0,\infty);L^{1}(\mathbb{R}^{m}))$。

现设 $u_{0}\in L^{\infty}(\mathbb{R}^{m})$。对 $r>0$，记 $\chi_{r}$ 为球 $\mathscr{B}_{r}(0)$ 的特征函数，$u^{r}$ 为\eqref{align:ch06-6-1-1}以 $\chi_{r}u_{0}\in L^{\infty}(\mathbb{R}^{m})\cap L^{1}(\mathbb{R}^{m})$ 为初值的可容许弱解。当 $r\to\infty$ 时，$\chi_{r}u_{0}\to u_{0}$（在 $L^{1}_{\mathrm{loc}}(\mathbb{R}^{m})$ 中）。因而由\eqref{align:ch06-6-2-9}，族 $\{u^{r}\}$ 在 $L^{1}_{\mathrm{loc}}$ 中收敛于某函数 $u$。显然 $u$ 是\eqref{align:ch06-6-1-1}、\eqref{align:ch06-6-1-2}的可容许弱解。由\cref{corollary:ch06-6-2-4}，此解唯一。由\cref{corollary:ch06-6-2-5}，当 $r$ 充分大时，在 $\mathbb{R}^{m}\times[0,\infty)$ 的任一紧子集上 $u\equiv u^{r}$。由于 $u^{r}(\cdot,t)\in C^{0}([0,\infty);L^{1}(\mathbb{R}^{m}))$，可知 $u(\cdot,t)\in C^{0}([0,\infty);L^{1}_{\mathrm{loc}}(\mathbb{R}^{m}))$。证明完成。
\end{proof}

\end{document}
